\documentclass[11pt]{article}

\usepackage[margin=1in]{geometry}
\usepackage{booktabs}
\usepackage{xcolor}
\usepackage{url}
\usepackage[colorlinks=true,linkcolor=blue!50!black,urlcolor=blue!50!black]{hyperref}

\title{Epidemiology Glossary}
\author{Centers for Disease Control and Prevention (CDC)\\
Reproductive Health}
\date{Page dated May 14, 2024; retrieved November 8, 2026}

\begin{document}
\maketitle

\noindent Source: \url{https://www.cdc.gov/reproductive-health/glossary/}

\noindent A glossary of commonly used terms in public health
surveillance and epidemiology, reproduced for offline reference.
Content of the U.S. Centers for Disease Control and Prevention.

\section*{A--B}

\begin{description}
\item[Age-adjusted mortality rate] A mortality rate statistically
modified to eliminate the effect of different age distributions
in the different populations.
\item[Agent] A factor, such as a microorganism, chemical
substance, or form of radiation, whose presence, excessive
presence, or (in deficiency diseases) relative absence is
essential for the occurrence of a disease.
\item[Age-specific mortality rate] A mortality rate limited to a
particular age group. The numerator is the number of deaths in
that age group; the denominator is the number of persons in that
age group in the population.
\item[Analytic epidemiology] The aspect of epidemiology concerned
with the search for health-related causes and effects. Uses
comparison groups, which provide baseline data, to quantify the
association between exposures and outcomes, and test hypotheses
about causal relationships.
\item[Analytic study] A comparative study intended to identify
and quantify associations, test hypotheses, and identify causes.
Two common types are cohort study and case-control study.
\item[Applied epidemiology] The application or practice of
epidemiology to address public health issues.
\item[Association] Statistical relationship between two or more
events, characteristics, or other variables.
\item[Attack rate] A variant of an incident rate, applied to a
narrowly defined population observed for a limited period of
time, such as during an epidemic.
\item[Attributable proportion] A measure of the public health
impact of a causative factor; proportion of a disease in a group
that is exposed to a particular factor which can be attributed to
their exposure to that factor.
\item[Bar chart] A visual display of the size of the different
categories of a variable. Each category or value of the variable
is represented by a bar.
\item[Bias] Deviation of results or inferences from the truth, or
processes leading to such systematic deviation. Any trend in the
collection, analysis, interpretation, publication, or review of
data that can lead to conclusions that are systematically
different from the truth.
\item[Biologic transmission] The indirect vector-borne
transmission of an infectious agent in which the agent undergoes
biologic changes within the vector before being transmitted to a
new host.
\item[Box plot] A visual display that summarizes data using a
``box and whiskers'' format to show the minimum and maximum
values (ends of the whiskers), interquartile range (length of
the box), and median (line through the box).
\end{description}

\section*{C}

\begin{description}
\item[Carrier] A person or animal without apparent disease who
harbors a specific infectious agent and is capable of
transmitting the agent to others. The carrier state may occur in
an individual with an infection that is inapparent throughout its
course (known as asymptomatic carrier), or during the incubation
period, convalescence, and postconvalescence of an individual
with a clinically recognizable disease. The carrier state may be
of short or long duration (transient carrier or chronic carrier).
\item[Case] In epidemiology, a countable instance in the
population or study group of a particular disease, health
disorder, or condition under investigation. Sometimes, an
individual with the particular disease.
\item[Case-control study] A type of observational analytic study.
Enrollment into the study is based on presence (``case'') or
absence (``control'') of disease. Characteristics such as
previous exposure are then compared between cases and controls.
\item[Case definition] A set of standard criteria for deciding
whether a person has a particular disease or health-related
condition, by specifying clinical criteria and limitations on
time, place, and person.
\item[Case-fatality rate] The proportion of persons with a
particular condition (cases) who die from that condition. The
denominator is the number of incident cases; the numerator is
the number of cause-specific deaths among those cases.
\item[Cause of disease] A factor (characteristic, behavior,
event, etc.) that directly influences the occurrence of disease.
A reduction of the factor in the population should lead to a
reduction in the occurrence of disease.
\item[Cause-specific mortality rate] The mortality rate from a
specified cause for a population. The numerator is the number of
deaths attributed to a specific cause during a specified time
interval; the denominator is the size of the population at the
midpoint of the time interval.
\item[Census] The enumeration of an entire population, usually
with details being recorded on residence, age, sex, occupation,
ethnic group, marital status, birth history, and relationship to
head of household.
\item[Chain of infection] A process that begins when an agent
leaves its reservoir or host through a portal of exit, and is
conveyed by some mode of transmission, then enters through an
appropriate portal of entry to infect a susceptible host.
\item[Class interval] A span of values of a continuous variable
which are grouped into a single category for a frequency
distribution of that variable.
\item[Cluster] An aggregation of cases of a disease or other
health-related condition, particularly cancer and birth defects,
which are closely grouped in time and place. The number of cases
may or may not exceed the expected number; frequently the
expected number is not known.
\item[Cohort] A well-defined group of people who have had a
common experience or exposure, who are then followed up for the
incidence of new diseases or events, as in a cohort or
prospective study. A group of people born during a particular
period or year is called a birth cohort.
\item[Cohort study] A type of observational analytic study.
Enrollment into the study is based on exposure characteristics
or membership in a group. Disease, death, or other
health-related outcomes are then ascertained and compared.
\item[Common source outbreak] An outbreak that results from a
group of persons being exposed to a common noxious influence,
such as an infectious agent or toxin. If the group is exposed
over a relatively brief period of time, so that all cases occur
within one incubation period, then the common source outbreak is
further classified as a point source outbreak. In some common
source outbreaks, persons may be exposed over a period of days,
weeks, or longer, with the exposure being either intermittent or
continuous.
\item[Confidence interval] A range of values for a variable of
interest, for instance, a rate, constructed so that this range
has a specified probability of including the true value of the
variable. The specified probability is called the confidence
level, and the end points of the confidence interval are called
the confidence limits.
\item[Confidence limit] The minimum or maximum value of a
confidence interval.
\item[Contact] Exposure to a source of an infection, or a person
so exposed.
\item[Contagious] Capable of being transmitted from one person to
another by contact or close proximity.
\item[Contingency table] A two-variable table with
cross-tabulated data.
\item[Control] In a case-control study, comparison group of
persons without disease.
\item[Crude mortality rate] The mortality rate from all causes of
death for a population.
\item[Cumulative frequency] In a frequency distribution, the
number or proportion of cases or events with a particular value
or in a particular class interval, plus the total number or
proportion of cases or events with smaller values of the
variable.
\item[Cumulative frequency curve] A plot of the cumulative
frequency rather than the actual frequency for each class
interval of a variable. This type of graph is useful for
identifying medians, quartiles, and other percentiles.
\end{description}

\section*{D}

\begin{description}
\item[Death-to-case ratio] The number of deaths attributed to a
particular disease during a specified time period divided by the
number of new cases of that disease identified during the same
time period.
\item[Demographic information] The person: characteristics -- age,
sex, race, and occupation -- of descriptive epidemiology used to
characterize the populations at risk.
\item[Denominator] The lower portion of a fraction used to
calculate a rate or ratio. In a rate, the denominator is usually
the population (or population experience, as in person-years,
etc.) at risk.
\item[Dependent variable] In a statistical analysis, the outcome
variable(s) or the variable(s) whose values are a function of
other variable(s) (called independent variable(s) in the
relationship under study).
\item[Descriptive epidemiology] The aspect of epidemiology
concerned with organizing and summarizing health-related data
according to time, place, and person.
\item[Determinant] Any factor, whether event, characteristic, or
other definable entity, that brings about change in a health
condition, or in other defined characteristics.
\item[Direct transmission] The immediate transfer of an agent
from a reservoir to a susceptible host by direct contact or
droplet spread.
\item[Distribution] In epidemiology, the frequency and pattern of
health-related characteristics and events in a population. In
statistics, the observed or theoretical frequency of values of a
variable.
\item[Dot plot] A visual display of the actual data points of a
noncontinuous variable.
\item[Droplet nuclei] The residue of dried droplets that may
remain suspended in the air for long periods, may be blown over
great distances, and are easily inhaled into the lungs and
exhaled.
\item[Droplet spread] The direct transmission of an infectious
agent from a reservoir to a susceptible host by spray with
relatively large, short-ranged aerosols produced by sneezing,
coughing, or talking.
\end{description}

\section*{E}

\begin{description}
\item[Endemic disease] The constant presence of a disease or
infectious agent within a given geographic area or population
group; may also refer to the usual prevalence of a given disease
within such area or group.
\item[Environmental factor] An extrinsic factor (geology,
climate, insects, sanitation, health services, etc.) which
affects the agent and the opportunity for exposure.
\item[Epidemic] The occurrence of more cases of disease than
expected in a given area or among a specific group of people over
a particular period of time.
\item[Epidemic curve] A histogram that shows the course of a
disease outbreak or epidemic by plotting the number of cases by
time of onset.
\item[Epidemic period] A time period when the number of cases of
disease reported is greater than expected.
\item[Epidemiologic triad] The traditional model of infectious
disease causation. Includes three components: an external agent,
a susceptible host, and an environment that brings the host and
agent together, so that disease occurs.
\item[Epidemiology] The study of the distribution and
determinants of health-related states or events in specified
populations, and the application of this study to the control of
health problems.
\item[Evaluation] A process that attempts to determine as
systematically and objectively as possible the relevance,
effectiveness, and impact of activities in the light of their
objectives.
\item[Experimental study] A study in which the investigator
specifies the exposure category for each person (clinical trial)
or community (community trial), then follows the person or
community to detect the effects of the exposure.
\item[Exposed (group)] A group whose members have been exposed to
a supposed cause of disease or health state of interest, or
possess a characteristic that is a determinant of the health
outcome of interest.
\end{description}

\section*{F--H}

\begin{description}
\item[Frequency distribution] A complete summary of the
frequencies of the values or categories of a variable; often
displayed in a two-column table: the left column lists the
individual values or categories, the right column indicates the
number of observations in each category.
\item[Frequency polygon] A graph of a frequency distribution with
values of the variable on the x-axis and the number of
observations on the y-axis; data points are plotted at the
midpoints of the intervals and are connected with a straight
line.
\item[Graph] A way to show quantitative data visually, using a
system of coordinates.
\item[Health] A state of complete physical, mental, and social
well-being and not merely the absence of disease or infirmity.
\item[Health indicator] A measure that reflects, or indicates,
the state of health of persons in a defined population, for
instance, the infant mortality rate.
\item[Health information system] A combination of health
statistics from various sources, used to derive information
about health status, health care, provision and use of services,
and impact on health.
\item[High-risk group] A group in the community with an elevated
risk of disease.
\item[Histogram] A graphic representation of the frequency
distribution of a continuous variable. Rectangles are drawn in
such a way that their bases lie on a linear scale representing
different intervals, and their heights are proportional to the
frequencies of the values within each of the intervals.
\item[Host] A person or other living organism that can be
infected by an infectious agent under natural conditions.
\item[Host factor] An intrinsic factor (age, race, sex,
behaviors, etc.) which influences a person's exposure,
susceptibility, or response to a causative agent.
\item[Hyperendemic disease] A disease that is constantly present
at a high incidence and/or prevalence rate.
\item[Hypothesis] A supposition, arrived at from observation or
reflection, that leads to refutable predictions. Any conjecture
cast in a form that will allow it to be tested and refuted.
\item[Hypothesis, null] The first step in testing for statistical
significance in which it is assumed that the exposure is not
related to disease.
\item[Hypothesis, alternative] The hypothesis, to be adopted if
the null hypothesis proves implausible, in which exposure is
associated with disease.
\end{description}

\section*{I--K}

\begin{description}
\item[Immunity, active] Resistance developed in response to
stimulus by an antigen (infecting agent or vaccine) and usually
characterized by the presence of antibody produced by the host.
\item[Immunity, herd] The resistance of a group to invasion and
spread of an infectious agent, based on the resistance to
infection of a high proportion of individual members of the
group. The resistance is a product of the number susceptible and
the probability that those who are susceptible will come into
contact with an infected person.
\item[Immunity, passive] Immunity conferred by an antibody
produced in another host and acquired naturally by an infant
from its mother or artificially by administration of an
antibody-containing preparation (antiserum or immune globulin).
\item[Incidence rate] A measure of the frequency with which an
event, such as a new case of illness, occurs in a population over
a period of time. The denominator is the population at risk; the
numerator is the number of new cases occurring during a given
time period.
\item[Incubation period] A period of subclinical or inapparent
pathologic changes following exposure, ending with the onset of
symptoms of infectious disease.
\item[Independent variable] An exposure, risk factor, or other
characteristic being observed or measured that is hypothesized
to influence an event or manifestation (the dependent variable).
\item[Indirect transmission] The transmission of an agent carried
from a reservoir to a susceptible host by suspended air particles
or by animate (vector) or inanimate (vehicle) intermediaries.
\item[Individual data] Data that have not been put into a
frequency distribution or rank ordered.
\item[Infectivity] The proportion of persons exposed to a
causative agent who become infected by an infectious disease.
\item[Inference, statistical] In statistics, the development of
generalizations from sample data, usually with calculated degrees
of uncertainty.
\item[Interquartile range] The central portion of a distribution,
calculated as the difference between the third quartile and the
first quartile; this range includes about one-half of the
observations in the set, leaving one-quarter of the observations
on each side.
\end{description}

\section*{L--M}

\begin{description}
\item[Latency period] A period of subclinical or inapparent
pathologic changes following exposure, ending with the onset of
symptoms of chronic disease.
\item[Mean, arithmetic] The measure of central location commonly
called the average. It is calculated by adding together all the
individual values in a group of measurements and dividing by the
number of values in the group.
\item[Mean, geometric] The mean or average of a set of data
measured on a logarithmic scale.
\item[Measure of association] A quantified relationship between
exposure and disease; includes relative risk, rate ratio, odds
ratio.
\item[Measure of central location] A central value that best
represents a distribution of data. Measures of central location
include the mean, median, and mode. Also called the measure of
central tendency.
\item[Measure of dispersion] A measure of the spread of a
distribution out from its central value. Measures of dispersion
used in epidemiology include the interquartile range, variance,
and the standard deviation.
\item[Median] The measure of central location which divides a set
of data into two equal parts.
\item[Medical surveillance] The monitoring of potentially exposed
people to detect early symptoms of disease.
\item[Midrange] The halfway point or midpoint in a set of
observations. For most types of data, it is calculated as the sum
of the smallest observation and the largest observation, divided
by two. For age data, one is added to the numerator. The midrange
is usually calculated as an intermediate step in determining
other measures.
\item[Mode] A measure of central location, the most frequently
occurring value in a set of observations.
\item[Morbidity] Any departure, subjective or objective, from a
state of physiological or psychological well-being.
\item[Mortality rate] A measure of the frequency of occurrence of
death in a defined population during a specified interval of
time.
\item[Mortality rate, infant] A ratio expressing the number of
deaths among children under 1 year of age reported during a given
time period divided by the number of births reported during the
same time period. The infant mortality rate is usually expressed
per 1,000 live births.
\item[Mortality rate, neonatal] A ratio expressing the number of
deaths among children from birth up to but not including 28 days
of age divided by the number of live births reported during the
same time period. The neonatal mortality rate is usually
expressed per 1,000 live births.
\item[Mortality rate, postneonatal] A ratio expressing the number
of deaths among children from 28 days up to but not including 1
year of age during a given time period divided by the number of
lives births reported during the same time period. The
postneonatal mortality rate is usually expressed per 1,000 live
births.
\end{description}

\section*{N--O}

\begin{description}
\item[Natural history of disease] The temporal course of disease
from onset (inception) to resolution.
\item[Necessary cause] A causal factor whose presence is required
for the occurrence of the effect (of disease).
\item[Nominal scale] Classification into unordered qualitative
categories; for instance, race, religion, and country of birth as
measurements of individual attributes are purely nominal scales,
as there is no inherent order to their categories.
\item[Normal curve] A bell-shaped curve that results when a
normal distribution is graphed.
\item[Normal distribution] The symmetrical clustering of values
around a central location. The properties of a normal
distribution include the following: (1) It is a continuous,
symmetrical distribution; both tails extend to infinity; (2) the
arithmetic mean, mode, and median are identical; and, (3) its
shape is completely determined by the mean and standard
deviation.
\item[Numerator] The upper portion of a fraction.
\item[Observational study] Epidemiological study in situations
where nature is allowed to take its course. Changes or
differences in one characteristic are studied in relation to
changes or differences in others, without the intervention of
the investigator.
\item[Odds ratio] A measure of association which quantifies the
relationship between an exposure and health outcome from a
comparative study; also known as the cross-product ratio.
\item[Ordinal scale] Classification into ordered qualitative
categories; for instance, social class (I, II, III, etc.), where
the values have a distinct order, but their categories are
qualitative in that there is no natural (numerical) distance
between their positive values.
\item[Outbreak] Synonymous with epidemic. Sometimes the preferred
word, as it may escape sensationalism associated with the word
epidemic. Alternatively, a localized as opposed to generalized
epidemic.
\end{description}

\section*{P}

\begin{description}
\item[Pandemic] An epidemic occurring over a very wide area
(several countries or continents) and usually affecting a large
proportion of the population.
\item[Pathogenicity] The proportion of persons infected, after
exposure to a causative agent, who then develop clinical disease.
\item[Percentile] The set of numbers from 0 to 100 that divide a
distribution into 100 parts of equal area, or divide a set of
ranked data into 100 class intervals with each interval
containing 1/100 of the observations. A particular percentile,
say the 5th percentile, is a cut point with 5\% of the
observations below it and the remaining 95\% of the observations
above it.
\item[Period prevalence] The amount a particular disease present
in a population over a period of time.
\item[Person-time rate] A measure of the incidence rate of an
event, for instance, a disease or death, in a population at risk
over an observed period to time, that directly incorporates time
into the denominator.
\item[Pie chart] A circular chart in which the size of each
``slice'' is proportional to the frequency of each category of a
variable.
\item[Point prevalence] The amount of a particular disease
present in a population at a single point in time.
\item[Population] The total number of inhabitants of a given area
or country. In sampling, the population may refer to the units
from which the sample is drawn, not necessarily the total
population of people.
\item[Predictive value positive] A measure of the predictive
value of a reported case or epidemic; the proportion of cases
reported by a surveillance system or classified by a case
definition which are true cases.
\item[Prevalence] The number or proportion of cases or events or
conditions in a given population.
\item[Prevalence rate] The proportion of persons in a population
who have a particular disease or attribute at a specified point
in time or over a specified period of time.
\item[Propagated outbreak] An outbreak that does not have a
common source, but instead spreads from person to person.
\item[Proportion] A type of ratio in which the numerator is
included in the denominator. The ratio of a part to the whole,
expressed as a ``decimal fraction'' (for instance, 0.2), as a
fraction (1/5), or, loosely, as a percentage (20\%).
\item[Proportionate mortality] The proportion of deaths in a
specified population over a period of time attributable to
different causes. Each cause is expressed as a percentage of all
deaths, and the sum of the causes must add to 100\%. These
proportions are not mortality rates, since the denominator is
all deaths, not the population in which the deaths occurred.
\item[Public health surveillance] The systematic collection,
analysis, interpretation, and dissemination of health data on an
ongoing basis, to gain knowledge of the pattern of disease
occurrence and potential in a community, in order to control and
prevent disease in the community.
\end{description}

\section*{Q--R}

\begin{description}
\item[Race-specific mortality rate] A mortality rate limited to a
specified racial group. Both numerator and denominator are
limited to the specified group.
\item[Random sample] A sample derived by selecting persons such
that each person has the same probability of selection.
\item[Range] In statistics, the difference between the largest
and smallest values in a distribution. In common use, the span of
values from smallest to largest.
\item[Rate] An expression of the frequency with which an event
occurs in a defined population.
\item[Rate ratio] A comparison of two groups in terms of
incidence rates, person-time rates, or mortality rates.
\item[Ratio] The value obtained by dividing one quantity by
another.
\item[Relative risk] A comparison of the risk of some
health-related event such as disease or death in two groups.
\item[Representative sample] A sample whose characteristics
correspond to those of the original population or reference
population.
\item[Reservoir] The habitat in which an infectious agent
normally lives, grows and multiplies; reservoirs include human
reservoirs, animals reservoirs, and environmental reservoirs.
\item[Risk] The probability that an event will occur, for
instance, that a person will become ill or die within a stated
period of time or age.
\item[Risk factor] An aspect of personal behavior or lifestyle,
an environmental exposure, or an inborn or inherited
characteristic that is associated with an increased occurrence
of disease or other health-related event or condition.
\item[Risk ratio] A comparison of the risk of some health-related
event such as disease or death in two groups.
\end{description}

\section*{S}

\begin{description}
\item[Sample] A selected subset of a population. A sample may be
random or non-random and it may be representative or
non-representative.
\item[Scatter diagram] A graph in which each dot represents
paired values for two continuous variables, with the x-axis
representing one variable and the y-axis representing the other;
used to display the relationship between the two variables; also
called a scattergram.
\item[Seasonality] Change in physiological status or in disease
occurrence that conforms to a regular seasonal pattern.
\item[Secondary attack rate] A measure of the frequency of new
cases of a disease among the contacts of known cases.
\item[Secular trend] Changes over a long period of time,
generally years or decades.
\item[Sensitivity] The ability of a system to detect epidemics
and other changes in disease occurrence. The proportion of
persons with disease who are correctly identified by a screening
test or case definition as having disease.
\item[Sentinel surveillance] A surveillance system in which a
pre-arranged sample of reporting sources agrees to report all
cases of one or more notifiable conditions.
\item[Sex-specific mortality rate] A mortality rate among either
males or females.
\item[Skewed] A distribution that is asymmetrical.
\item[Specificity] The proportion of persons without disease who
are correctly identified by a screening test or case definition
as not having disease.
\item[Sporadic] A disease that occurs infrequently and
irregularly.
\item[Spot map] A map that indicates the location of each case of
a rare disease or outbreak by a place that is potentially
relevant to the health event being investigated, such as where
each case lived or worked.
\item[Standard deviation] The most widely used measure of
dispersion of a frequency distribution, equal to the positive
square root of the variance.
\item[Standard error (of the mean)] The standard deviation of a
theoretical distribution of sample means about the true
population mean.
\item[Sufficient cause] A causal factor or collection of factors
whose presence is always followed by the occurrence of the
effect (of disease).
\item[Surveillance] See Public health surveillance.
\item[Survival curve] A curve that starts at 100\% of the study
population and shows the percentage of the population still
surviving at successive times for as long as information is
available. May be applied not only to survival as such, but also
to the persistence of freedom from a disease, or complication or
some other endpoint.
\end{description}

\section*{T--V}

\begin{description}
\item[Table] A set of data arranged in rows and columns.
\item[Table shell] A table that is complete except for the data.
\item[Transmission of infection] Any mode or mechanism by which
an infectious agent is spread through the environment or to
another person.
\item[Trend] A long-term movement or change in frequency, usually
upwards or downwards.
\item[Universal precautions] Recommendations issued by CDC to
minimize the risk of transmission of bloodborne pathogens,
particularly HIV (human immunodeficiency virus) and HBV
(hepatitis B virus), by health care and public safety workers.
Barrier precautions are to be used to prevent exposure to blood
and certain body fluids of all patients.
\item[Validity] The degree to which a measurement actually
measures or detects what it is supposed to measure.
\item[Variable] Any characteristic or attribute that can be
measured.
\item[Variance] A measure of the dispersion shown by a set of
observations, defined by the sum of the squares of deviations
from the mean, divided by the number of degrees of freedom in
the set of observations.
\item[Vector] An animate intermediary in the indirect
transmission of an agent that carries the agent from a reservoir
to a susceptible host.
\item[Vehicle] An inanimate intermediary in the indirect
transmission of an agent that carries the agent from a reservoir
to a susceptible host.
\item[Virulence] The proportion of persons with clinical disease
who, after becoming infected, become severely ill or die.
\item[Vital statistics] Systematically tabulated information
about births, marriages, divorces, and deaths, based on
registration of these vital events.
\end{description}

\section*{W--Z}

\begin{description}
\item[Years of potential life lost] A measure of the impact of
premature mortality on a population, calculated as the sum of the
differences between some predetermined minimum or desired life
span and the age of death for individuals who died earlier than
that predetermined age.
\item[Zoonoses] An infectious disease that is transmissible under
normal conditions from animals to humans.
\end{description}

\end{document}
